\chapter*{Introduction}
\addcontentsline{toc}{chapter}{Introduction}

Optimal control problems with nonconvex cost functions are widely used in a great number of practical and theorical problems, such as inverse problems in biomedical or seismic imaging, certain energy efficiency models (see \cite{vogel}), image restoration (see \cite{hintermueller}), compressed sensing (see \cite{foucart}), optimal control problems (see \cite{casas}), and other applications. Their importance relies on the properties that the nonconvexity can provide to the optimal control solution, such as the sparsity. Of course, this desirable feature over the control is not only induced exclusively by nonconvex cost functions; for example, Georg Stadler \cite{stadler2006} proposed a control problem with a $L^1$ penalizer governned by a elliptic partial differential equation and control pointwise constrains, and showed that the optimal control solution is sparse over a subset of its domain. Therefore, we see that it is possible to induce sparsity over the optimal control of a problem choosing a suitable penalization. However, the cost functional used in \cite{stadler2006} is convex but not differentiable.\\

Thus, the natural question which arises from the sparsity of the control is that if we can `maximise' as much as possible the subset of its domain where the sparsity appears. One possibility to address this is the use a $L^0$ penalizer in the cost function \cite{ito-kunisch2014}, because this one is linked to the minimization of the measure of the support of the control. However, this penalization is not even a norm-type penalization, and lack of convexity and differentiability is evident. Hence, in order to induce sparsity of the optimal control we can use this kind of penalizations at cost of the loss of convexity of the cost function. In fact, the work of Kasumufi Ito and Karl Kunisch \cite{ito-kunisch2014} show us that, if we use a $L^q$ penalizer, with $q\in [0,1)$, sparsity of the control is assured. Unfortunately, the loss of convexity is inherent to this type of problems, causing that the standard tools of the Convex Analysis and Functional Analysis can not be used for proving aspects such as existence and unicity of solutions for the problem and its regularity, or deducing optimally systems and subsequently the numerical treatment of this problems. \\

Between the $L^1$ and $L^0$ penalizers, it can be considerated $L^q$ penalizers with $q\in [0,1)$, which, in some sense, can be considered as natural aproximations of the $L^0$ penalizer \cite{merino2019} as $q\to 0$. However, there are notable features of the sparsity induced in the optimal control which arise from this particular kind of penalized problems, specially that the jump discontinuities of the optimal control can be characterized by a constant with depends on the problem parameters and appears at its boundary's support. This property can be used in practical problems which needs a sharp localized action of the control within the domain. Despite this, the lack of convexity in the cost function is inevitable, which makes  its theoretical and numerical treatment challenging as we pointed out before.\\

In order to study the sparsity induced by the $L^q$ penalizer, a regularization or a relaxation of this term can be used, see for example \cite{ekeland}, \cite{shane}, \cite{warga}. However, it could happen that the these properties become weaker than the non regularized problem. Instead of this, Ito and Kunisch in \cite{ito-kunisch2014} propose a Pontryagin-type principle with allow us to characterise the jumps of the control under certain hypothesis of the functions involved in the control problem. One advantage of this approach is that, if a control problem satisfies the hypothesis proposed in \cite{ito-kunisch2014}, and has an optimal control, then its sparsity can be studied through a uni dimensional problem. However, the imposition of the existence of a solution for the control problem is required, being the real existence of it as an assumption. In fact, the optimal control problem treated in \cite{ito-kunisch2014} involves an $H_{0}^{1}$ norm of the control, which allow them to prove existence of an optimal solution under compacity arguments. Pedro Merino in \cite{merino2019} study a similar problem than Ito and Kunisch, but changing the $H_0^1$ norm with a classical $L^2$ Tikhonov term. Nevertheless, a Huber-type regularization is proposed in order to carry out with the nonconvex $L^q$ term, making that the sparse properties of the control dependents on the Huber regularization term. Despite this, he proved that the regularized control problem has an optimal solution, and moreover, a optimally system can be deduced with enough regularity to propose a numerical scheme to solve it, using the so called DC-Algorithm. A latter work of the same author propose a Semi Smooth Newton method for solving the problem.\\

The question of the existence of an optimal solution for this kind of problems, without using regularization or the non convex term such as additional terms in the cost function that allow the use of compacity arguments was thereby open. This problem motivate us to investigate a way to obtain the existence of an optimal control solution for this kind of problems, and later to characterise its sparsity properties. The use of the tools proposed by Ito and Kunisch are useful for our work, and allow us to extend them for deriving an optimally system that allow us to construct a solution for the $L^q$ problem, without regularizing or relaxing it.\\

On the other hand, it can be proposed a relaxation of the cost function of the problem which involves both the Tikhonov term and the nonconvex penalty $L^q$, using the so called  \textit{Convex Regularization} of the uni dimensional version of this terms. Specifically, we discover that the convexified problem allow us to construct an optimal solution for the original problem, without the requirement of the regularization of the non convex term using additional parameters such as in the Huber regularization case, and only through its optimally system. Moreover, we find that the optimally system of the convexified problem is nearly equivalent than the first approach we present, but each one constructed by different theoretical approaches. In addition, the convexified problem provide us tools to deduce a priori error estimates for solving numerically the original problem using a finite elements approach.\\

The optimal control problems discussed so far may have point-wise control restrictions. In fact, the Pontryagin-type principle used in our investigation is derived for a control problems with this kind of restrictions; however, if we impose point-wise state constraints instead of control ones in the non convex penalized problems, we deal with additional difficulties in order to prove the existence of solutions, such as the derivation of an optimally system that allow us to study numerically the problem. Optimal control problems with point-wise state constrains but convex cost functions are widely studied, we refer to \cite{casas, casas-clason-kunisch2012, casas-delosreyes}. The main difficult present in this kind of problems is that their optimally systems present Lagrange multipliers with low regularity, such as Borel regular measures. In order to deal with this issue, Christian Meyer et al. proposed in \cite{meyer-2007} a Lavrentiev-type regularization with mix state and control in a new artificial control variable, which depends on a regularization parameter. This change of variable allow them to derive a more regular optimally system which is used to derive a numerical method for solving the problem. However, the convexity of the cost function is crucial for studying the limit of the family of problems linked to the regularization parameter.\\

In our particular case, the lack of convexity difficult the theoretical treatment of the problem with point-wise state constrains. However, the convex relaxation, plus a combination of mixed control-state point-wise constrains in the problem allow us to study the pure state constrained one. \\

Respect to the numerical treatment of our problem, we present three different schemes: the Sequential Quadratic Hamiltonian Algorithm, the Difference of Convex functions Algorithm, and the Semismooth Newton Algorithm. The first scheme can be applied directly to the control constrained problem without a regularization or a explicit computation of its optimally system. However, in order to ensure its convergence, a sufficient descent condition which depends on the Hamiltonian of the problem must be imposed, and therefore, this method could be restrictive for a general optimal control problems. In addition, since this method needs to solve a point-wise one dimensional minimization problem at each node of the discretized domain, an additional computational power could be needed. The first issue about this method can be carried out by the favorable structure of the Hamiltonian of our problem, while the second one can be solved with parallelization of the computations.\\

The second method, Difference of Convex algorithm, was first introduced in a finite dimensional optimization frame. Several works about this method are \cite{dinh, dihn2, dihn3, dinh4}, and many other papers. In contrast, the application of this method in infinite dimensional context is novel, and works such as of Yongdo Lim et. al \cite{yongdo} present extensions to the DC algortihm for optimizaction in Hilber spaces. In the particular case of control problems with $L^q$ penalizers, \cite{merino2019} presents a numerical scheme  to the regularized control problem based on the DC algorithm. The main advantage behind this method is that, since the cost function can be split into the difference of two convex functions, and one of them has a closed subdifferential, it is easy to construct part of the algorithm explicit. Despite this, to end a complete iteration of the algorithm, a $L^1$ control constrained problem must be solved.\\

The third method, Semi smooth Newton algorithm, was proposed for solving system of equations which have certain regularity (semismoothness regularity). The finite dimensional form of this algorithm was introduced by Mifflin in \cite{mifflin}, and its extension for infinite dimensional spaces was introduced by several wors, included \cite{qi, qisun}. For the $L^q$ penalized control constrained problem, \cite{merino2021} proposed an hybrid method which combines the Semismooth Newton method and the DC method. However, this method is specially useful for the numerical treatment of the $L^q$ penalized method with point-wise state constrains, but applied first with a regularization first over the constrains, following a Lavrentiev type regularization, proposed by \cite{meyer-2007}.\\

Our contributions in this work are the following:

\begin{itemize}
    \item We prove the existence of a generalized control solution to a $L^q$ penalized problem with control point-wise constrains, defining a thresholding operator which extends the methodology proposed by \cite{ito-kunisch2014}.
    \item  We prove the existence of a solution for the penalized $L^q$ control problem, with mixed control-state point-wise constrains, such as a pure point-wise state constrain problem.
    \item We establish the best numerical algorithm for solving the penalized problem with control point-wise constrains, facing up three different numerical algorithms, through a extensively set of numerical experiments.
    \item We address a model...
\end{itemize}

This work is organized as follows: the first chapter is devoted to a review of some theoretical results about Functional Analysis and Convex Analysis that will be useful through the thesis. 

